\section[Conception et accompagnement]{Concevoir les services et accompagner les usagers dans un système en mutation}
Aujourd'hui, la mise en oeuvre d'une bibliothèque numérique ne peut se limiter uniquement à l'assemblage d'un répertoire de fichiers statiques. Pour que la nouvelle plateforme s'impose comme un véritable outil scientifique, l'établissement doit mettre en place une forte médiation documentaire. Pour cela, la bibliothèque Cujas doit transformer un stock de fichiers en un environnement de connaissances éditorialisées. Aujourd'hui, le rôle des bibliothèques ne se limite pas à stocker des fichiers isolés; il s'agit d'assembler, de contextualiser et de faire dialoguer ces ressources entre elles pour en faire émerger une plus value intellectuelle et de nouvelles connaissances. Dans leur ouvrage, Antoine Isaac, Emmanuelle Bermès et Gautier Poupeau expliquent ce phénomène à travers le concept de la \enquote{citation de contenus}. Pour les auteurs, elle consiste à \enquote{sélectionner des ressources sur le Web, de les réorganiser, de les annoter et de donner à consulter ce travail dans une interface particulière permettant de saisir le sens qui émerge de la juxtaposition de ces ressources}. Dans ce troisième âge, ce traitement intellectuel est l'un des moyens utilisés par lese institutions patrimoniales, transformant le document-objet en un espace d'apprentissage et de valorisation du travail scientifique réalisés par des spécialistes. 

Le premier jalon de cette attente réside dans la réintroduction du paratexte d'accompagnement autour des oeuvres numérisées. Les constats issus des enquêtes de terrain à la BIU Cujas sont univoques: les usagers soulignent l'importance de mettre en valeur ce qui est numérisé avec un texte descriptif. Cette attente est vigoureusement relayée par les chercheurs interrogés. L'un des professeurs insiste sur le fait qu'il est capital \enquote{de présenter les ouvrages et d'expliciter leur qualité pour susciter la curiosité scientifique de publics}, permettant aux spécialistes et au grand public de s'y intéresser notamment si on lie cette éditorialisation à des notions d'actualité.

L'une des pistes envisagées, serait la création d'une \enquote{galerie des grands juristes} sous forme d'expositions virtuelles et de parcours thématiques commentés. Laurent Pfister propose \enquote{d'associer directement la communauté de recherche}ce projet d'éditorialisation, \enquote{en confiant à des doctorants ou à des jeunes chercheurs la rédaction de courts articles scientifiques}. Ces articles expliciteraient une thématique ou un sujet de recherche où le matériel utilisé fut les documents numérisés par la bibliothèque Cujas. Cette éditorialisation partagée présente un double bénéfice : elle valorise les travaux des jeunes chercheurs tout en tissant un lien entre la bibliothèque et les laboratoires de recherche.

L'éditorialisation est possible sous différentes formes en fonctions des formats ou des protocoles utilisés par la bibliothèque, façonnant l'espace numérique. Comme le souligne Elina Leblanc, elle s'appuie sur la définition de Marcello Vitali-Rosati en disant que \enquote{les protocoles employés pour numériser et partager des contenus dans un environnement numérique contribuent à définir plusieurs manières de percevoir le patrimoine  écrit numérique}. L'éditorialisation, la création de parcours thématiques guidés et l'ajout de notices d'autorité vulgarisées sont devenus des services indispensables pour accompagner l'étudiant dans l'appropriation des sources complexes. En mettant en place ce type de contenu, cela permet de répondre à l'une des attentes d'une doctorant concernant le découragement étudiant face à la masse de documents numérisés. Elina Leblanc préconise une médiation active afin de réduire la surcharge cognitive et la désorientation des lecteurs.
\bigskip

La future bibliothèque numérique de la bibliothèque Cujas doit devenir un outil de référence pour la recherche juridique, offrant un environnement où la rigueur académique est soutenue et non entravée, par la technologie. C'est à travers les suggestions faites par les chercheurs eux-mêmes, qui permettront selon eux à la bibliothèque de gagner en rigueur scientifique. Face à ces exigences, le logiciel XTF ne pouvait plus répondre à leurs attentes. C'est pourquoi le prototype Codex, développé sous Python-Django, a été pensé pour implémenter techniquement les standards issus du Web de données, notamment via des API et le protocole IIIF.